\chapter[Equipo de radiografía industrial]{Equipo para la radiografía industrial con fuente radiactiva}

\section{Irriditrón modelo 520}

Esta unidad de radiación gamma está diseñada para operar en forma remota y puede manejar hasta 100 Ci de iridio-192, con lo que se reducen los tiempos de exposición a unos cuantos minutos. El iridio-192 es uno de los radionúclidos más versátiles y puede aplicarse a secciones delgadas o gruesas, de 6.25 a 76 mm de espesor para el acero o su equivalente para otros materiales.

El equipo básico consta de:
\begin{itemize}
  \item cabezal blindado;
  \item manipulador con indicador y cable de control;
  \item ducto para la fuente radiactiva.
\end{itemize}

Puede tener los siguientes accesorios:
\begin{itemize}
  \item colimador de tungsteno, con cono de $70^\circ$;
  \item conjunto de abrazadera y trípode;
  \item ducto para fuente más corto.
\end{itemize}

\section{Cabezal blindado}

El cabezal blindado (figura~\ref{fig:cabezal-blindado}) consiste en un recipiente cerrado de acero inoxidable que contiene suficiente material aislante (uranio agotado) para reducir el nivel de radiación exterior. El blindaje está diseñado para aceptar con seguridad una fuente de 100 Ci de iridio-192, centrada en un tubo de titanio en forma de ``S''.

\begin{figure}[htbp]
  \centering
  \includegraphics[width=0.72\textwidth]{cabezalBlindado.png}
  \caption{Cabezal blindado y sus componentes.}
  \label{fig:cabezal-blindado}
\end{figure}

Los componentes indicados en la figura son:
\begin{description}
  \item[a)] Tapón con rosca para el extremo del ducto.
  \item[b)] Tapón con rosca para el cable de control.
  \item[c)] Tubo ``S''.
  \item[d)] Fuente.
  \item[e)] Agarradera.
  \item[f)] Chapa.
  \item[g)] Placa indicadora de material radiactivo.
\end{description}

En el exterior del cabezal existe un mecanismo de seguridad que impide el movimiento de la fuente en el tubo ``S'', para evitar que pueda salir accidentalmente durante la preparación. El cabezal pesa 18 kg y está fabricado para cumplir con todos los requisitos de diseño del Departamento de Transportes de F. U.

\section{Manipulador con indicador y cable de control}

\begin{figure}[htbp]
  \centering
  \includegraphics[width=0.86\textwidth]{manipulador.png}
  \caption{Manipulador con indicador y cable de control.}
  \label{fig:manipulador}
\end{figure}

Este conjunto (figura~\ref{fig:manipulador}) proporciona los medios para accionar la fuente desde el cabezal y, a través del ducto, llevarla a la posición de exposición. Se cuenta con 15.2 m de cable de acero inoxidable, impulsado por un engranaje especial dentro de dos conductores exteriores. En un extremo del cable impulsor se encuentra un gancho que se conecta al cable de la fuente para su operación.

Al girar la manija gira el engranaje, moviendo de esta manera el cable impulsor y la fuente fuera del cabezal. En la sección de la manija está montado un medidor que proporciona una lectura de la distancia recorrida por la fuente a través del ducto. Este conjunto manipulador trasladará la fuente, fuera del cabezal, una distancia de 7 m.

\section{Ducto para la fuente radiactiva}

\begin{figure}[htbp]
  \centering
  \includegraphics[width=0.86\textwidth]{ducto.png}
  \caption{Ducto para la fuente radiactiva.}
  \label{fig:ducto-fuente}
\end{figure}

El ducto estándar tiene 7 m de longitud y sirve como guía para trasladar la fuente desde el cabezal hasta el punto en el cual se va a tomar la radiografía. El ducto consiste en una manguera tejida, flexible, a prueba de intemperie y con accesorios de acero inoxidable en sus extremos.

Un extremo cuenta con una cuerda apropiada para enroscarse y fijarse al cabezal. El otro extremo cuenta con una tapa de cabeza de muñón que limita el recorrido de la fuente e impide que esta salga. La tapa de cabeza de muñón es removible, lo que permite las siguientes operaciones:
\begin{enumerate}[label=\alph*)]
  \item conectar en el extremo del ducto el colimador de tungsteno;
  \item conectar la tapa de cabeza de muñón directamente a la salida del cabezal para tomar radiografías panorámicas dentro de un cuarto de radiografía o detrás de una barrera blindada.
\end{enumerate}

\section{Accesorios}

\begin{enumerate}
  \item \textbf{Colimador de tungsteno.} Es una pieza que puede acoplarse directamente a la salida del cabezal cuando se desea tomar radiografías con haz direccional. También puede acoplarse al extremo del ducto de la fuente para tomar radiografías con haz direccional y con el cabezal alejado de la fuente. Tiene un peso de 4.2 kg y cuenta con una rosca de acero inoxidable.

  \item \textbf{Conjunto de abrazadera y trípode.} Consiste en un trípode universal con una abrazadera especial para el extremo del ducto de la fuente. Esto permite al operador un ajuste continuo y completo de la elevación y de la posición rotacional de la tapa de cabeza de muñón, para radiografías panorámicas, o la colocación exacta del colimador de tungsteno.

  \item \textbf{Ducto para fuente más corto.} Tiene una longitud de 3 m y es similar al ducto para fuente estándar, excepto que no incluye la tapa de cabeza de muñón. Este ducto sirve para:
  \begin{itemize}
    \item usarse como extensión del ducto para fuente estándar, con lo que se logra una longitud total de 10 m;
    \item utilizarse en lugar del ducto estándar cuando las condiciones de blindaje lo permitan. La tapa de cabeza de muñón y el colimador de tungsteno pueden acoplarse al extremo de este ducto.
  \end{itemize}
\end{enumerate}

\section{Instrucciones de operación}

Para unir el ducto de la fuente al cabezal, se quita el tapón con rosca (a en la figura~\ref{fig:cabezal-blindado}) y en su lugar se atornilla el acople del extremo correspondiente del ducto. El cabezal debe estar bloqueado.

Para conectar el manipulador con indicador, se quita el tapón con rosca (b en la figura~\ref{fig:cabezal-blindado}) y se conecta el cable de control al cable de la fuente, teniendo mucho cuidado para que no se zafe.

Para disponer el equipo para la operación, el extremo del ducto de la fuente se coloca en la posición donde se va a tomar la radiografía. Su fijación puede hacerse mediante un trípode, cinta adhesiva, abrazadera, bloque de madera con perforaciones u otra forma adecuada. El ducto debe estar tan recto como sea posible para facilitar la operación.

Una vez realizado lo anterior, se extiende el cable de control a toda su longitud y se coloca la manija de control tan alejada como sea posible del extremo del ducto de la fuente donde quedará expuesto el iridio. El operador también debe estar tan alejado como sea posible del punto donde se va a utilizar la fuente, preferentemente detrás de un blindaje para el personal.

Después de abrir el Irriditrón con la llave proporcionada, se expone la fuente girando la manija de control hasta que esta se detenga. El indicador de posición muestra en todo momento la distancia recorrida por la fuente. No se debe forzar la manija.

Una vez realizada la exposición, se regresa la fuente girando la manija en sentido contrario, hasta que el indicador marque «cero». Cierre la chapa de seguridad y retire la llave para que el equipo no sea utilizado por personal no autorizado.

\section{Seguridad en el manejo de materiales radiactivos}

Es conveniente utilizar un monitor de radiación durante la operación, para tener la seguridad de que la fuente está blindada de forma apropiada y verificar que no se excedan los límites de exposición.

Por ningún motivo se permite trasladar el equipo ensamblado para hacer nuevas exposiciones en tramos largos o cambiar de nivel con el equipo ensamblado. Se deberá desensamblar y mantener los equipos de monitoreo encendidos hasta que la fuente quede completamente blindada dentro del contenedor.

Al término de cada jornada de trabajo, el personal ocupacionalmente expuesto deberá asegurarse de que la fuente quede bien blindada dentro del contenedor, revisando los niveles de radiación que presenta y comparándolos con los que tenía antes de iniciar la jornada.

Al término de la jornada, el material deberá guardarse en el Almacén de Material Radiactivo de la compañía. Cuando por necesidades de servicio la unidad radiográfica se encuentre fuera de la ciudad donde se localiza el almacén de material radiactivo, deberá mantenerse en el portacontenedor de la unidad y guardarse en un estacionamiento cerrado y con vigilancia. Por ningún motivo se deberá dejar la unidad estacionada en la vía pública con el material radiactivo en ella y sin vigilancia.

\section{Almacenamiento y transporte}

El almacenamiento y transporte del material radiactivo dentro del territorio nacional se rige por las disposiciones y normas establecidas por la CNSNS (Comisión Nacional de Seguridad Nuclear y Salvaguardias).

A la unidad en que se transporte material radiactivo se le debe colocar un rótulo o etiqueta en cada uno de sus lados. En la mitad superior irá el símbolo internacional que indica la presencia de radiación, de color negro, sobre un fondo contrastante. En la mitad inferior del rótulo se indicará la leyenda: «NO PERMANEZCA INNECESARIAMENTE CERCA DE ESTE VEHÍCULO».

Además, el vehículo deberá estar provisto de un cordón amarillo, necesario para acordonar dos veces el perímetro de la unidad, postes para colocar el cordón y letreros con el símbolo internacional que indica la presencia de radiación.

El tránsito de un vehículo que contenga material radiactivo se hará por la ruta más corta, siempre observando los límites de velocidad señalados por las autoridades.

Únicamente personal autorizado (POE) de la empresa podrá transportar material radiactivo en las unidades debidamente equipadas y destinadas para ello. Cuando se transporte material radiactivo por necesidades del servicio a otra ciudad, se deberá seguir la ruta más corta y segura, y reportarse como mínimo una vez al día a las oficinas generales, indicando el lugar donde se encuentra y la ruta que está siguiendo.

Si por alguna falla mecánica se tiene que detener la unidad, nunca y por ningún motivo se dejará sola en la vía pública con el material radiactivo. El auxiliar o el técnico permanecerá al cuidado de la unidad y una persona deberá conseguir o pedir ayuda mecánica. La unidad deberá ser revisada mecánicamente antes de efectuar un viaje de más de 100 km; además, deberá contar con llanta de refacción, gato hidráulico, cables pasa-corriente, letreros preventivos, cordón con una longitud de dos veces el diámetro de la camioneta para acordonarla en caso necesario, letreros preventivos de peligro radiactivo, extintor de fuego y una caja de herramientas básicas para alguna falla mecánica.

En caso de que la unidad donde se transporte el material radiactivo sufra alguna falla y no pueda ser reparada en el lugar, por ningún motivo se transportará la fuente en compartimentos ocupados por pasajeros (autobús, ferrocarril, avión o correo). Se deberá llamar a las oficinas generales para el envío de otra unidad equipada para el transporte.

La intensidad máxima permitida fuera del vehículo y en la cabina es de $2\ \mathrm{mR/h}$.

La unidad nunca se podrá dejar sola en la vía pública cuando las fuentes se encuentren dentro de ella.

Cuando el material radiactivo no esté en uso, las fuentes se deben guardar en un depósito destinado a este fin. El lugar de almacenamiento debe estar adecuadamente protegido y debidamente señalizado. Deberá asignarse a una persona competente que lo tenga a su cargo.

Únicamente el personal autorizado puede introducir o retirar material radiactivo de dicho local. Este lugar debe estar protegido contra toda intrusión. Los lugares de almacenamiento deben elegirse de modo que el riesgo de inundación sea mínimo y deben inspeccionarse periódicamente para detectar una posible contaminación.

Se deberá llevar un registro de todas las fuentes radiactivas que se encuentren almacenadas.
